%==============================================================================
% Evaluatie
%==============================================================================

\chapter{Evaluatie}
\label{ch:evaluatie}

Dit hoofdstuk evalueert de voorgestelde ERP-migratie op basis van de uitgewerkte TO-BE-situatie en de uitgevoerde Proof of Concept. De evaluatie focust op procesverbetering, operationele efficiëntie, documenttraceerbaarheid, voorraadactualiteit en voorbereiding op Peppol-readiness.

Omdat de oorspronkelijke werkomgeving van Hanuman BV beperkt historische logdata bevatte, is de evaluatie gebaseerd op observaties tijdens de AS-IS-analyse, indicatieve tijdsinschattingen, scenario-gebaseerde simulaties en validatie binnen de gecontroleerde Odoo-testomgeving. De vermelde KPI-cijfers moeten daarom geïnterpreteerd worden als indicatieve vergelijkingswaarden en niet als harde productiemetingen.

De evaluatie heeft niet als doel om een volledige productie-implementatie te vervangen, maar om na te gaan of de voorgestelde gefaseerde migratie technisch en organisatorisch haalbaar is voor een servicegerichte Belgische KMO.

\section{KPI-evaluatiekader}
\label{sec:kpi-evaluatiekader}

Om de impact van de voorgestelde ERP-implementatie gestructureerd te kunnen beoordelen, werd een evaluatiekader opgesteld op basis van Key Performance Indicators (KPI's).

Deze KPI's werden geselecteerd op basis van:
\begin{itemize}
    \item de geïdentificeerde knelpunten in de AS-IS-analyse;
    \item de doelstellingen van de implementatie;
    \item de onderzoeksvraag van deze bachelorproef;
    \item relevante procesindicatoren uit de literatuur rond ERP-implementatie.
\end{itemize}

De evaluatie focust voornamelijk op:
\begin{itemize}
    \item reductie van manuele handelingen;
    \item verbetering van procesdoorlooptijd;
    \item vermindering van foutgevoeligheid;
    \item verbetering van traceerbaarheid;
    \item voorraadactualiteit;
    \item voorbereiding op elektronische facturatie.
\end{itemize}

\subsection{Totstandkoming van de KPI-baseline}

De AS-IS-baseline werd niet opgebouwd uit automatisch geregistreerde proceslogs, omdat Hanuman BV in de oorspronkelijke situatie geen geïntegreerd systeem gebruikte dat dergelijke metingen structureel vastlegde. De baseline werd daarom indicatief bepaald op basis van drie bronnen: shadowing gedurende één werkdag, analyse van bestaande documenten en informele validatie via stakeholdergesprekken.

De TO-BE-waarden zijn gebaseerd op de uitgevoerde PoC-scenario's in Odoo. Daarbij werd nagegaan hoeveel stappen binnen de gekozen testscenario's nodig waren en hoe Odoo de documentflow, voorraadmutaties, facturatie en traceerbaarheid ondersteunde. De vergelijking tussen AS-IS en TO-BE is dus scenario-gebaseerd en bedoeld om de verwachte richting en grootteorde van procesverbetering te onderbouwen.

De cijfers in dit hoofdstuk mogen daarom niet gelezen worden als statistisch bewezen productieresultaten. Ze geven een onderbouwde indicatie van de mogelijke procesimpact binnen de gecontroleerde PoC-context.

\subsection{Geselecteerde KPI's}

Tabel~\ref{tab:kpis} geeft een overzicht van de gebruikte KPI's.

\begin{table}[H]
\centering
\caption{Gebruikte KPI's voor evaluatie}
\label{tab:kpis}
\begin{tabular}{|p{4cm}|p{5cm}|p{4cm}|}
\hline
\textbf{KPI} & \textbf{Beschrijving} & \textbf{Meetmethode} \\
\hline
Aantal manuele stappen & Aantal handmatige acties binnen een procesflow & Observatie en scenarioanalyse \\
\hline
Doorlooptijd & Tijd tussen start en afronding van een proces & Indicatieve tijdsinschatting per scenario \\
\hline
Correctiemomenten & Aantal noodzakelijke manuele correcties & Analyse van procesflows en observaties \\
\hline
Documenttraceerbaarheid & Mogelijkheid om gekoppelde documenten op te volgen & Procesvalidatie in Odoo \\
\hline
Voorraadactualiteit & Mate waarin voorraadwijzigingen tijdig verwerkt worden & Simulatie van stockmutaties \\
\hline
Peppol-readiness & Voorbereiding op elektronische facturatie & Configuratiecontrole in testomgeving \\
\hline
\end{tabular}
\end{table}

Naast de kwalitatieve beoordeling werd per kernproces ook een indicatieve vergelijking gemaakt tussen de AS-IS-situatie en de TO-BE-situatie. Tabel~\ref{tab:kpi-resultaten} vat de belangrijkste KPI-resultaten samen.

\begin{table}[H]
\centering
\caption{Indicatieve, scenario-gebaseerde vergelijking van KPI-resultaten}
\label{tab:kpi-resultaten}
\begin{tabular}{|p{4cm}|p{3cm}|p{3cm}|p{3cm}|}
\hline
\textbf{KPI} & \textbf{Indicatieve AS-IS} & \textbf{TO-BE in PoC} & \textbf{Indicatief effect} \\
\hline
Manuele stappen verkoopflow & ca. 8 & ca. 4 & Mogelijke daling van ongeveer 50\% \\
\hline
Doorlooptijd verkoopflow & ca. 15--20 min & ca. 5--8 min & Indicatieve daling van ongeveer 55--65\% \\
\hline
Correctiemomenten per 10 factuurflows & ca. 2--3 & ca. 0--1 & Sterke daling verwacht \\
\hline
Zoektijd servicedossier & ca. 5--10 min & ca. 1--2 min & Sterke daling verwacht \\
\hline
Voorraadactualiteit & Laattijdig bijgewerkt & Bij validatie bijgewerkt & Sterke verbetering verwacht \\
\hline
Documenttraceerbaarheid & Beperkt & Geïntegreerd & Sterke verbetering \\
\hline
Peppol-readiness & Beperkt & Voorbereid binnen testomgeving & Sterke verbetering binnen scope \\
\hline
\end{tabular}
\end{table}

De cijfers tonen aan dat vooral de verkoop- en facturatieflow efficiënter kan verlopen binnen de gekozen PoC-scenario's. De grootste winst ontstaat door het hergebruik van centrale klant- en productgegevens en door de automatische koppeling tussen offerte, verkooporder, levering en factuur.

Naast de samenvattende KPI-resultaten werd per kernproces een meer gedetailleerde vergelijking gemaakt. Deze vergelijking is gebaseerd op de processtappen uit de AS-IS-analyse en de overeenkomstige standaardflows binnen Odoo, zoals offerte naar verkooporder, verkooporder naar levering, verkooporder naar factuur en repair naar onderdelenverbruik en facturatie.

\begin{table}[H]
\centering
\caption{Indicatieve KPI-vergelijking per kernproces}
\label{tab:kpi-detailprocessen}
\begin{tabular}{|p{4cm}|p{3cm}|p{3cm}|p{3cm}|}
\hline
\textbf{Proces / KPI} & \textbf{Indicatieve AS-IS} & \textbf{TO-BE in PoC} & \textbf{Effect} \\
\hline
Offerte opstellen & ca. 8--12 min & ca. 3--5 min & Sneller door hergebruik klant- en productdata \\
\hline
Offerte naar factuur & ca. 15--20 min & ca. 5--8 min & Minder overnamewerk door automatische documentflow \\
\hline
Aantal manuele stappen verkoopflow & ca. 8 & ca. 4 & Mogelijke daling van ongeveer 50\% \\
\hline
Klantgegevens opnieuw controleren & ca. 3--4 keer & ca. 1 keer & Centrale contactfiche \\
\hline
Voorraadupdate na verkoop of herstelling & Vaak einde werkdag & Bij validatie van transactie & Actueler voorraadbeeld \\
\hline
Zoektijd servicedossier & ca. 5--10 min & ca. 1--2 min & Centrale historiek per klant/toestel \\
\hline
Service naar facturatie & ca. 20--30 min & ca. 10--15 min & Arbeid en onderdelen sneller factureerbaar \\
\hline
Onderdelenregistratie bij herstelling & Manueel & Gekoppeld aan voorraadmutatie & Minder risico op vergeten onderdelen \\
\hline
Correctiemomenten per 10 factuurflows & ca. 2--3 & ca. 0--1 & Minder foutgevoelige overname \\
\hline
Peppol-controlepunten & Verspreid / manueel & Centraal controleerbaar & Betere voorbereiding op e-invoicing \\
\hline
\end{tabular}
\end{table}

De gedetailleerde vergelijking toont dat de verwachte efficiëntiewinst niet voortkomt uit één afzonderlijke automatisatie, maar uit de combinatie van centrale stamdata, automatische documentkoppeling en geïntegreerde voorraadregistratie. Vooral processen waarbij informatie in de AS-IS-situatie meerdere keren opnieuw gebruikt of gecontroleerd moest worden, profiteren van de standaard Odoo-flow.

\section{Vergelijking tussen AS-IS en TO-BE situatie}
\label{sec:asis-tobe}

Op basis van de uitgewerkte scenario's werd een vergelijking gemaakt tussen de huidige situatie en de voorgestelde TO-BE-situatie binnen Odoo.

\subsection{Verkoop- en facturatieproces}

In de huidige situatie bestaat het verkoopproces uit meerdere afzonderlijke stappen waarbij gegevens manueel moeten worden opgezocht, gecontroleerd en opnieuw ingevoerd.

Binnen de TO-BE-situatie worden offertes, verkooporders, leveringen en facturen automatisch aan elkaar gekoppeld binnen dezelfde ERP-omgeving. Scenario 1 valideerde dit aan de hand van verkooporder S00001 voor Cafe Central BV, waarbij de flow van offerte naar levering, factuur en betalingsopvolging werd uitgevoerd. Scenario 4 valideerde daarnaast een gemengde order met een onderhoudsdienst en fysieke producten.

Tabel~\ref{tab:salesvergelijking} toont de vergelijking tussen beide situaties.

\begin{table}[H]
\centering
\caption{Indicatieve vergelijking verkoop- en facturatieproces}
\label{tab:salesvergelijking}
\begin{tabular}{|p{5cm}|p{3cm}|p{3cm}|}
\hline
\textbf{Indicator} & \textbf{Indicatieve AS-IS} & \textbf{TO-BE in PoC} \\
\hline
Aantal manuele stappen & ca. 8 & ca. 4 \\
\hline
Gemiddelde doorlooptijd & ca. 15--20 minuten & ca. 5--8 minuten \\
\hline
Correctiemomenten per 10 factuurflows & ca. 2--3 & ca. 0--1 \\
\hline
Documentkoppeling & Beperkt & Geïntegreerd \\
\hline
Dubbele gegevensinvoer & Frequent & Beperkt \\
\hline
\end{tabular}
\end{table}

De resultaten tonen aan dat vooral de automatische koppeling tussen offerte, verkooporder en factuur een belangrijke impact heeft op de efficiëntie van het proces. In de oorspronkelijke situatie moesten klantgegevens, productinformatie en factuurgegevens op meerdere momenten opnieuw worden opgezocht of ingevoerd. Binnen de TO-BE-situatie worden deze gegevens hergebruikt vanuit centrale stamdata.

Daarnaast vermindert de centrale beschikbaarheid van klant- en productgegevens het risico op foutieve of inconsistente informatie. Hierdoor daalt niet enkel de verwerkingstijd, maar ook de kans op correcties achteraf.

\subsection{Voorraadbeheer}

Binnen de oorspronkelijke situatie werden voorraadwijzigingen vaak pas op het einde van de werkdag administratief verwerkt. Hierdoor ontstonden tijdelijke verschillen tussen administratieve en werkelijke voorraad.

In de TO-BE-situatie worden voorraadbewegingen gekoppeld aan verkoop-, aankoop- en repairprocessen. Dit werd in de PoC gevalideerd via drie scenario's. In scenario 2 steeg de voorraad van Pomp 15 bar van 3 naar 8 na ontvangst van aankooporder P00001. In scenario 3 daalde dezelfde voorraad van 8 naar 7 na verbruik in repair order WH/RO/00001. In scenario 4 daalde de voorraad van Waterfilter cartridge van 36 naar 34 na levering aan OfficeBeans NV.

\begin{table}[H]
\centering
\caption{Vergelijking voorraadbeheer}
\label{tab:voorraadvergelijking}
\begin{tabular}{|p{5cm}|p{3cm}|p{3cm}|}
\hline
\textbf{Indicator} & \textbf{AS-IS} & \textbf{TO-BE in PoC} \\
\hline
Actueel voorraadbeeld & Beperkt & Beschikbaar na validatie \\
\hline
Automatische stockmutaties & Nee & Ja, binnen uitgevoerde scenario's \\
\hline
Moment voorraadupdate & Vaak einde werkdag & Bij validatie van transactie \\
\hline
Risico op voorraadverschillen & Hoog & Beperkt mits correcte registratie \\
\hline
Koppeling met verkoop & Beperkt & Geïntegreerd \\
\hline
Koppeling met herstellingen & Manueel & Gekoppeld aan repair flow \\
\hline
\end{tabular}
\end{table}

De verwerking van stockmutaties zorgt voor een actueler voorraadbeeld en verlaagt het risico op foutieve voorraadinschattingen. Dit is vooral relevant voor onderdelen die gebruikt worden bij herstellingen, omdat een laattijdige registratie rechtstreeks kan leiden tot vertragingen of verkeerde bestellingen.

\subsection{Service-, herstel- en onderhoudsprocessen}

Binnen de oorspronkelijke situatie werd service-informatie verspreid bijgehouden via e-mails, facturen en manuele notities.

De TO-BE-situatie laat toe om klantgegevens, toestellen, onderdelen, uitgevoerde werken en onderhoudsopdrachten centraler te koppelen binnen één ERP-omgeving. Scenario 3 valideerde de repair flow voor Hotel Aroma BV, inclusief onderdelenverbruik en facturatie. Scenario 5 valideerde serienummertracking voor Espresso Pro X met serienummer SN-EPX-2026-002 en een equipment record EPX-2026-002 met preventieve onderhoudsplanning.

\begin{table}[H]
\centering
\caption{Indicatieve vergelijking service-, herstel- en onderhoudsprocessen}
\label{tab:servicevergelijking}
\begin{tabular}{|p{5cm}|p{3cm}|p{3cm}|}
\hline
\textbf{Indicator} & \textbf{Indicatieve AS-IS} & \textbf{TO-BE in PoC} \\
\hline
Centrale historiek & Beperkt & Beschikbaar per klant/toestel \\
\hline
Traceerbaarheid onderdelen & Manueel & Gekoppeld aan repair en voorraad \\
\hline
Koppeling met facturatie & Beperkt & Geïntegreerd \\
\hline
Opvolging service-interventies & Verspreid & Centraler opgevolgd \\
\hline
Serienummertracking & Beperkt of manueel & Gevalideerd via SN-EPX-2026-002 \\
\hline
Preventief onderhoud & Manueel of ad hoc & Gepland via equipment record \\
\hline
Zoektijd historiek & ca. 5--10 minuten & ca. 1--2 minuten \\
\hline
Service naar facturatie & ca. 20--30 minuten & ca. 10--15 minuten \\
\hline
\end{tabular}
\end{table}

De integratie van service-, voorraad- en onderhoudsprocessen verlaagt het risico op vergeten onderdelen en onvolledige facturatie. Bovendien maakt centrale historiek het eenvoudiger om terugkerende problemen bij een toestel of klant op te volgen.

\section{Evaluatie van procesverbeteringen}
\label{sec:procesverbeteringen}

Op basis van de uitgewerkte scenario's kan worden vastgesteld dat de voorgestelde ERP-implementatie verschillende operationele verbeteringen kan opleveren.

\subsection{Reductie van manuele handelingen}

Een van de belangrijkste verbeteringen is de reductie van manuele gegevensinvoer. Door klantgegevens, producten en documenten centraal te beheren, moeten gegevens minder vaak opnieuw worden ingevoerd of gecontroleerd.

Binnen de scenario's werd indicatief vastgesteld dat het aantal manuele stappen binnen de standaard verkoopflow ongeveer kan halveren. Dit resultaat ligt boven de vooropgestelde doelstelling van 20--30\% reductie in manuele handelingen, maar moet gelezen worden als scenario-gebaseerde indicatie en niet als bewezen productiecijfer.

De reductie komt vooral voort uit:
\begin{itemize}
    \item hergebruik van klantgegevens;
    \item hergebruik van product- en prijsinformatie;
    \item automatische omzetting van offerte naar verkooporder;
    \item automatische koppeling tussen verkooporder en factuur;
    \item geïntegreerde documenthistoriek.
\end{itemize}

De reductie betekent niet dat menselijke controle volledig verdwijnt. In een KMO-context blijft controle door de gebruiker belangrijk, bijvoorbeeld bij prijsafspraken, uitzonderlijke servicekosten of onvolledige klantgegevens. De winst ligt vooral in het vermijden van herhaald overtypen en het beperken van administratieve tussenstappen.

\subsection{Verbeterde documenttraceerbaarheid}

Binnen de TO-BE-situatie worden offertes, verkooporders, leveringen, facturen en voorraadbewegingen aan elkaar gekoppeld.

Hierdoor wordt het eenvoudiger om:
\begin{itemize}
    \item openstaande processen op te volgen;
    \item historiek te raadplegen;
    \item fouten sneller te identificeren;
    \item documenten te reconstrueren.
\end{itemize}

Deze verbeterde traceerbaarheid werd in de PoC zichtbaar via de smart buttons en gekoppelde documenten bij verkooporders, leveringen en facturen. Dit is vooral relevant binnen servicegerichte processen waarbij meerdere interventies en onderdelen gekoppeld zijn aan dezelfde klant of hetzelfde toestel.

\subsection{Snellere facturatie}

De omzetting van verkooporders naar facturen vermindert de administratieve verwerkingstijd binnen de gekozen scenario's. Waar het opstellen en verzenden van een factuur in de AS-IS-situatie indicatief 15 tot 20 minuten in beslag nam, kon dit in de TO-BE-situatie binnen de PoC worden teruggebracht tot ongeveer 5 tot 8 minuten.

Daarnaast verlaagt de integratie van servicegegevens en onderdelen het risico op ontbrekende factuurlijnen of vergeten prestaties. Dit is vooral belangrijk bij herstellingen, waar facturatie vaak bestaat uit een combinatie van arbeid, onderdelen en eventuele verplaatsingskosten.

Bij servicefacturatie is de winst iets beperkter dan bij standaardverkoop, omdat diagnose, arbeidsregistratie en onderdelenkeuze nog steeds inhoudelijke beoordeling vragen. Toch kan de doorlooptijd ook hier dalen doordat onderdelen en prestaties minder vaak manueel moeten worden overgenomen naar een afzonderlijk facturatiedocument.

\subsection{Betere voorraadcontrole}

Doordat voorraadbewegingen gekoppeld worden aan transacties, ontstaat een actueler voorraadbeeld.

Hierdoor wordt:
\begin{itemize}
    \item het risico op stockverschillen beperkt;
    \item onderdelenverbruik beter geregistreerd;
    \item de afhankelijkheid van manuele controle verminderd;
    \item toekomstige aankoopbehoeften sneller zichtbaar.
\end{itemize}

Voor Hanuman BV is dit vooral relevant omdat onderdelenverbruik rechtstreeks samenhangt met service- en herstelactiviteiten. Een correct voorraadbeeld verkleint het risico dat herstellingen vertraging oplopen door ontbrekende onderdelen.

\section{Evaluatie van Peppol-readiness}
\label{sec:peppol-evaluatie}

Een belangrijk onderdeel van deze bachelorproef is de voorbereiding op gestructureerde elektronische facturatie via Peppol.

Tijdens de evaluatie werd nagegaan of de voorgestelde configuratie voldoende voorbereid is om elektronische facturen te kunnen genereren binnen de testomgeving. De PoC testte geen effectieve productie-uitwisseling via een gecertificeerd Peppol Access Point.

Daarbij werd gecontroleerd of:
\begin{itemize}
    \item btw-nummers correct geregistreerd zijn;
    \item verplichte adresgegevens aanwezig zijn;
    \item bankgegevens correct geconfigureerd zijn;
    \item factuurstructuren voldoende gestandaardiseerd zijn;
    \item stamdata centraal beheerd wordt;
    \item btw-tarieven en fiscale posities correct zijn ingesteld;
    \item Peppol in demomodus zichtbaar is;
    \item een UBL BIS 3 XML-bijlage zichtbaar is in de verzendwizard.
\end{itemize}

Tabel~\ref{tab:peppol-readiness} geeft een overzicht van de Peppol-readinesscontrole.

\begin{table}[H]
\centering
\caption{Peppol-readinesscontrole binnen de Odoo-testomgeving}
\label{tab:peppol-readiness}
\begin{tabular}{|p{5cm}|p{3cm}|p{5cm}|}
\hline
\textbf{Controlepunt} & \textbf{Status} & \textbf{Opmerking} \\
\hline
Btw-nummer onderneming & In orde & BE0477472701 geregistreerd \\
\hline
Adresgegevens onderneming & In orde & Belgisch adres aanwezig \\
\hline
Bankgegevens onderneming & In orde & KBC-testbankrekening aanwezig \\
\hline
Belgische btw-codes & In orde & Belgische fiscale lokalisatie beschikbaar \\
\hline
Klantgegevens Cafe Central BV & In orde & Belgisch adres en btw-gegevens aanwezig \\
\hline
Peppol-gerelateerde klantinformatie & In orde binnen PoC & Peppol-informatie zichtbaar op klantfiche \\
\hline
Factuurvelden & In orde & Facturen bevatten klant, lijnen, btw, totalen en betaalreferentie \\
\hline
Peppol Demo & In orde binnen testomgeving & Optie ``met Peppol (Demo)'' zichtbaar \\
\hline
UBL BIS 3 XML-bijlage & In orde binnen testomgeving & XML-bijlage zichtbaar in verzendwizard \\
\hline
Productie-uitwisseling via Access Point & Buiten scope & Niet uitgevoerd \\
\hline
\end{tabular}
\end{table}

Uit de controle blijkt dat de onderzochte Peppol-readinesspunten binnen de testomgeving in orde zijn. Het onderdeel dat bewust niet werd gevalideerd, is de effectieve productie-uitwisseling met een gecertificeerd Peppol Access Point. Dit valt buiten de scope van de Proof of Concept, maar vormt wel een noodzakelijke stap bij een latere productie-uitrol.

Op basis van deze controle kan gesteld worden dat de voorgestelde configuratie een bruikbare basis biedt voor verdere Peppol-voorbereiding. De centrale registratie van klant- en bedrijfsgegevens vermindert het risico op validatiefouten in vergelijking met de oorspronkelijke situatie.

Hoewel de Proof of Concept geen volledige productie-integratie met een Peppol Access Point omvat, toont de evaluatie aan dat Hanuman BV binnen de testomgeving technisch beter voorbereid is op toekomstige e-invoicingverplichtingen.

\section{Aansluiting tussen BPMN-modellen en PoC-scenario's}
\label{sec:bpmn-poc-aansluiting}

De BPMN-achtige AS-IS- en TO-BE-modellen uit Hoofdstuk~\ref{ch:implementatie-analyse} sluiten aan op de uitgevoerde PoC-scenario's. De modellen tonen op procesniveau waar manuele overdrachtsmomenten ontstaan en hoe deze in de TO-BE-situatie worden verminderd door centrale stamdata en gekoppelde documentflows.

De verkoop- en facturatieflow uit de TO-BE-modellen werd gevalideerd in scenario 1 en scenario 4. Scenario 1 toont de volledige flow van offerte naar verkooporder, levering, factuur en betalingsopvolging. Scenario 4 toont bijkomend dat een onderhoudsdienst en fysieke producten samen op één factuur kunnen worden verwerkt, terwijl enkel de fysieke producten een voorraadlevering veroorzaken.

Het TO-BE-voorraadproces werd gevalideerd in scenario 2, scenario 3 en scenario 4. Scenario 2 toont een voorraadverhoging na ontvangst van een aankooporder. Scenario 3 toont voorraadvermindering door onderdelenverbruik in een repair order. Scenario 4 toont voorraadvermindering na levering van fysieke producten.

Het service- en herstelproces werd gevalideerd in scenario 3 en scenario 5. Scenario 3 toont de repair flow met onderdelenverbruik en facturatie. Scenario 5 breidt deze logica uit met serienummertracking, equipment records en preventieve onderhoudsplanning.

Hierdoor sluiten de BPMN-modellen voldoende aan op de uitgevoerde PoC. Een aandachtspunt blijft dat de BPMN-modellen conceptueel zijn opgesteld en geen volledige BPMN 2.0-specificatie vormen. Ze dienen vooral om de overgang van AS-IS naar TO-BE visueel te verduidelijken.

\section{Beperkingen van het onderzoek}
\label{sec:beperkingen}

Hoewel de resultaten van de evaluatie positief zijn, kent dit onderzoek verschillende beperkingen.

Ten eerste werd gewerkt binnen een gecontroleerde Proof of Concept-omgeving en niet binnen een volledige productieomgeving. Hierdoor konden niet alle praktische omstandigheden van een reële implementatie worden nagebootst.

Daarnaast is de evaluatie gedeeltelijk gebaseerd op scenario-gebaseerde simulaties en observaties tijdens de AS-IS-analyse. Omdat historische logdata beperkt beschikbaar was, konden bepaalde KPI's slechts indicatief worden geëvalueerd. De percentages en tijdsinschattingen in dit hoofdstuk geven daarom een onderbouwde indicatie van procesverbetering, maar vormen geen statistisch bewezen productiemeting.

Ook langdurige gebruikersacceptatie en organisatorische veranderingen konden binnen de beperkte onderzoeksperiode niet volledig worden onderzocht.

Verder werd bewust gekozen om zoveel mogelijk gebruik te maken van standaardfunctionaliteiten binnen Odoo. Hierdoor werden geavanceerde maatwerkfunctionaliteiten buiten scope gehouden.

Tot slot zijn de resultaten voornamelijk toepasbaar op kleine servicegerichte ondernemingen met gelijkaardige processen en beperkte organisatorische complexiteit.

\section{Samenvattende evaluatie}
\label{sec:samenvattende-evaluatie}

De evaluatie toont aan dat een gefaseerde implementatie van Odoo een duidelijke meerwaarde kan bieden voor een servicegerichte Belgische KMO zoals Hanuman BV.

De voorgestelde ERP-omgeving laat toe om:
\begin{itemize}
    \item manuele administratie te verminderen;
    \item documenttraceerbaarheid te verbeteren;
    \item voorraadbeheer beter te integreren;
    \item service-, herstel- en onderhoudsprocessen centraler op te volgen;
    \item voorbereiding op Peppol-readiness te versterken.
\end{itemize}

Daarnaast toont de evaluatie aan dat een gefaseerde aanpak geschikt is binnen een KMO-context, omdat deze implementatierisico's beperkt en gebruikers toelaat om stapsgewijs vertrouwd te raken met de nieuwe processen.

De resultaten ondersteunen daarmee de centrale onderzoeksvraag van deze bachelorproef, met de nuance dat de vastgestelde verbeteringen gebaseerd zijn op een gecontroleerde PoC en indicatieve KPI-vergelijking. De evaluatie vormt zo de basis voor de conclusies in het volgende hoofdstuk.